

\documentclass[12pt, twoside]{report}

\usepackage[
  top=1in, bottom=1in,
  inner=1.25in, outer=1in,
  includeheadfoot
]{geometry}
\usepackage{fontenc}
\usepackage{inputenc}
\usepackage{lmodern}
\usepackage{microtype}
\usepackage{xcolor}
\usepackage{graphicx}
\usepackage{hyperref}
\usepackage{fancyhdr}
\usepackage{titlesec}
\usepackage{tcolorbox}
\usepackage{enumitem}
\usepackage{booktabs}
\usepackage{longtable}
\usepackage{array}
\usepackage{parskip}
\usepackage{float}
\usepackage{caption}

\definecolor{MacSyncPrimary}{HTML}{6B1A35}   % Maroon — chapter titles
\definecolor{NoteGray}{HTML}{F2F2F2}         % Light gray — note boxes
\definecolor{WarnGray}{HTML}{EBEBEB}         % Slightly darker gray — warnings
\definecolor{BorderGray}{HTML}{999999}       % Gray border
\definecolor{DarkGray}{HTML}{2B2B2B}         % Near-black text

\hypersetup{
  colorlinks=true,
  linkcolor=MacSyncPrimary,
  urlcolor=MacSyncPrimary,
  pdftitle={MacSync Staff and Admin Guide},
  pdfauthor={MacSync Team},
  bookmarksnumbered=true,
  bookmarksopen=true,
}

\titleformat{\chapter}[display]
  {\normalfont\huge\bfseries\color{MacSyncPrimary}}
  {\chaptertitlename\ \thechapter}{10pt}{\Huge}

\titleformat{\section}
  {\normalfont\Large\bfseries\color{black}}
  {\thesection}{1em}{}

\titleformat{\subsection}
  {\normalfont\large\bfseries\color{black}}
  {\thesubsection}{1em}{}

\pagestyle{fancy}
\fancyhf{}
\fancyhead[LE]{\small\leftmark}
\fancyhead[RO]{\small MacSync Staff \& Admin Guide}
\fancyfoot[LE,RO]{\thepage}
\fancyfoot[C]{\small\textcolor{gray}{McMaster Engineering Society}}
\renewcommand{\headrulewidth}{0.4pt}
\renewcommand{\footrulewidth}{0.4pt}

\tcbuselibrary{skins, breakable}

\newtcolorbox{notebox}{
  colback=NoteGray, colframe=BorderGray,
  fonttitle=\bfseries\small, title=Note,
  breakable, left=6pt, right=6pt, top=4pt, bottom=4pt
}

\newtcolorbox{warnbox}{
  colback=WarnGray, colframe=DarkGray,
  fonttitle=\bfseries\small, title=Warning,
  breakable, left=6pt, right=6pt, top=4pt, bottom=4pt
}

\newtcolorbox{tipbox}{
  colback=NoteGray, colframe=BorderGray,
  fonttitle=\bfseries\small, title=Tip,
  breakable, left=6pt, right=6pt, top=4pt, bottom=4pt
}

\newcommand{\uielem}[1]{\textbf{#1}}
\newcommand{\version}{1.0.0}
\newcommand{\releasedate}{April 2026}

%% ============================================================
%%  DOCUMENT START
%% ============================================================
\begin{document}

\begin{titlepage}
  \centering
  \vspace*{2cm}

  % \includegraphics[width=0.35\textwidth]{logo.png}\\[1.5cm]

  {\Huge\bfseries\color{MacSyncPrimary} MacSync}\\[0.4cm]
  {\Large\color{DarkGray} MES Large Events Management Platform}\\[0.2cm]
  \noindent\rule{\textwidth}{1pt}\\[0.4cm]
  {\LARGE\bfseries Staff \& Admin Guide}\\[0.2cm]
  {\large Version \version\ ---\ \releasedate}\\[2cm]

  {\large Mahad Ahmed \textbullet\ Abyan Jaigirdar \textbullet\ Prerna Prabhu\\
  Farhan Rahman \textbullet\ Ali Zia}\\[0.4cm]
  {\normalsize\textit{Capstone 2025--2026 | Streamliners}}\\[0.6cm]
  {\large McMaster University \textbullet\ McMaster Engineering Society}\\[1cm]

  \vfill
  {\small\textcolor{gray}{\copyright\ 2026 MacSync. All rights reserved.\\
  McMaster Engineering Society}}
\end{titlepage}

\pagenumbering{roman}
\setcounter{page}{1}

\chapter*{Document History}
\addcontentsline{toc}{chapter}{Document History}
\begin{longtable}{@{} p{2cm} p{2.5cm} p{4cm} p{5cm} @{}}
  \toprule
  \textbf{Version} & \textbf{Date} & \textbf{Author} & \textbf{Changes} \\
  \midrule
  \endhead
  1.0.0 & April 2026 & MacSync Team & Initial release \\
  \bottomrule
\end{longtable}

\tableofcontents
\clearpage

\pagenumbering{arabic}
\setcounter{page}{1}


\chapter{Introduction}

\section{About MacSync}

MacSync is an event management platform built for the McMaster Engineering Society (MES).
It enables MES clubs and staff to create, manage, and promote large-scale events.
Staff manage everything through the web-based admin portal at
\texttt{admin.macsync.org}, while members use the mobile app to discover and
attend events. Both interfaces share the same account system.

\section{Who This Guide Is For}

This guide is for \textbf{staff and administrators} --- anyone with a staff role
who uses the admin portal to create or manage events, as well as System Admins
who manage users, roles, and platform settings. For the member-facing mobile
app, refer to the \textit{MacSync Member Guide}.

\section{How to Use This Guide}

Start with Chapter~\ref{chap:started} to access the admin portal, then refer
to the relevant chapters for each feature. Callout boxes highlight important
information:

\begin{notebox}
Notes provide extra context or background information.
\end{notebox}

\begin{warnbox}
Warnings flag actions that could have unintended consequences.
\end{warnbox}

\begin{tipbox}
Tips share best practices and useful shortcuts.
\end{tipbox}

\section{System Requirements}
\begin{itemize}
  \item Modern web browser: Chrome, Firefox, Safari, or Edge.
  \item Active internet connection.
  \item A valid \texttt{@mcmaster.ca} account with a staff role or
        System Admin privileges.
\end{itemize}


\chapter{Getting Started}
\label{chap:started}

\section{The Two MacSync Portals}

MacSync has two interfaces that share the same account system:

\begin{table}[H]
  \centering
  \caption{MacSync Portals Overview}
  \begin{tabular}{@{} p{3.2cm} p{4cm} p{5.5cm} @{}}
    \toprule
    \textbf{Portal} & \textbf{Who It's For} & \textbf{What You Can Do} \\
    \midrule
    \textbf{Mobile App}
      & All users (Members and Staff)
      & Browse events, register, purchase tickets, manage seat selections,
        view notifications. Staff can also create and manage events. \\
    \addlinespace
    \textbf{Web Admin Portal}
      & Staff and System Admins only
      & Create and manage events, manage users and roles,
        view attendee details, download venue reports, send notifications. \\
    \bottomrule
  \end{tabular}
\end{table}

\begin{notebox}
Staff can use \textbf{both} portals. Use the mobile app as an attendee
and the web admin portal to manage events.
\end{notebox}

\section{Accessing the Admin Portal}

Navigate to \texttt{admin.macsync.org} in any modern web browser.

\begin{notebox}
Only \texttt{@mcmaster.ca} accounts with an Admin role, a custom staff role,
or System Admin privileges can access this portal. Regular Members will see
an \textit{``Access denied. Sign in with a system administrator or event
staff account''} message.
\end{notebox}

\section{Logging In}

\begin{enumerate}
  \item Go to \texttt{admin.macsync.org}.
  \item Enter your \texttt{@mcmaster.ca} email and password.
  \item Click \uielem{Sign in}.
\end{enumerate}

\section{Navigating the Admin Portal}

Once logged in, the sidebar provides access to:

\begin{itemize}
  \item \uielem{Dashboard} --- Overview of platform activity.
  \item \uielem{Events} --- Create, edit, delete, and view events you manage.
  \item \uielem{Manage Roles} --- Create and delete custom roles
        (System Admins only).
  \item \uielem{Users} --- View and manage all registered users
        (System Admins only).
  \item \uielem{Notifications} --- Send announcements to event attendees.
\end{itemize}


\chapter{User Roles \& Permissions}
\label{chap:roles}

\section{Role Overview}

MacSync uses a flexible role-based access control (RBAC) system. Every registered
user automatically receives the \textbf{Member} role. Additional roles can be
created by a System Admin and assigned to users to grant event management privileges.
A user can hold \textbf{multiple roles simultaneously}.

There are two built-in roles on every MacSync instance:

\begin{table}[H]
  \centering
  \caption{Built-in Roles}
  \begin{tabular}{@{} p{2.5cm} p{1.8cm} p{8.3cm} @{}}
    \toprule
    \textbf{Role} & \textbf{Protected} & \textbf{Description} \\
    \midrule
    Member  & Yes & Default role for every registered user. Can browse
                    and register for events. Cannot create or manage events. \\
    Admin   & No  & Staff role. Can create events and manage events
                    where their role matches the event's Managing Role. \\
    \bottomrule
  \end{tabular}
\end{table}

Custom roles (e.g.\ \texttt{MES\_ADMIN}) can be created from the
\uielem{Manage Roles} page and assigned to any combination of users.

\section{System Admin}

The \textbf{System Admin} privilege is a special toggle separate from roles,
found under \uielem{Advanced Settings} on any user's profile.

\begin{itemize}
  \item A System Admin has \textbf{unrestricted access} to all platform
        features, events, users, and settings with no role restrictions.
  \item Only an existing System Admin can grant or revoke this privilege.
  \item In production, this should be held by MES Presidents and senior executives.
\end{itemize}

\begin{warnbox}
System Admin is listed under \uielem{Advanced Settings --- Dangerous actions
and system-level changes}. Grant this privilege carefully --- it overrides
all role and event-level restrictions.
\end{warnbox}

\section{How the Managing Role Works}

Every event has a \textbf{Managing Role} field that determines which staff
users can edit, delete, and view attendee data for that event.

\begin{itemize}
  \item When creating an event, the managing role defaults to the creator's
        highest assigned role.
  \item Any user who holds the matching role can administer that event.
  \item Example: if an event's managing role is \texttt{MES\_ADMIN},
        all users with the \texttt{MES\_ADMIN} role can manage it.
  \item System Admins can manage all events regardless of managing role.
\end{itemize}

\begin{tipbox}
To restrict an event to a specific team, create a custom role
(e.g.\ \texttt{HATCH\_EXEC}), assign it only to that team's members,
and set it as the managing role when creating the event.
\end{tipbox}

\section{Creating a New Role (System Admin Only)}

\begin{enumerate}
  \item Navigate to \uielem{Manage Roles} in the sidebar.
  \item Under \uielem{Create a new role}, enter a role name.
  \item Click \uielem{Create role}.
\end{enumerate}

\begin{notebox}
Role names are case-sensitive. Use consistent naming conventions
(e.g.\ all caps with underscores: \texttt{CLUB\_EXEC}).
\end{notebox}

\section{Deleting a Role (System Admin Only)}

\begin{enumerate}
  \item Navigate to \uielem{Manage Roles}.
  \item Find the role under \uielem{Existing roles}.
  \item Click \uielem{Delete role}.
\end{enumerate}

\begin{warnbox}
Deleting a role removes it from all users currently assigned that role and
from any events using it as their managing role. The \textbf{Member} role
is protected and cannot be deleted.
\end{warnbox}

\section{Assigning Roles to a User (System Admin Only)}

\begin{enumerate}
  \item Navigate to \uielem{Users} and click the target user.
  \item Click \uielem{Show advanced settings}.
  \item Under \uielem{Roles}, check the boxes for all roles to assign.
  \item Click \uielem{Save roles}.
\end{enumerate}

The user's active roles are displayed in the top-right of their profile page.

\section{Granting System Admin Privileges (System Admin Only)}

\begin{enumerate}
  \item Navigate to \uielem{Users} and open the target user's profile.
  \item Click \uielem{Show advanced settings}.
  \item Toggle \uielem{System admin} to the on position.
\end{enumerate}

\begin{warnbox}
There is no confirmation dialog for this action. Verify you have the correct
user's profile open before enabling System Admin.
\end{warnbox}


\chapter{Managing Events}

\section{Creating an Event}

Any user with an Admin role or a custom staff role can create events.
Regular Members cannot create events.

\begin{enumerate}
  \item From the sidebar, click \uielem{Events}, then click
        \uielem{Create event} in the top-right corner.
  \item Fill in the event details (see Section~\ref{sec:eventfields}).
  \item Configure optional signup settings if needed
        (see Sections~\ref{sec:table} and~\ref{sec:bus}).
  \item Set the \uielem{Managing role} to control which staff can
        administer this event.
  \item Click \uielem{Create event} to publish.
\end{enumerate}

\section{Event Form Fields}
\label{sec:eventfields}

\begin{table}[H]
  \centering
  \caption{Create Event --- Field Reference}
  \begin{tabular}{@{} p{3cm} p{1.5cm} p{8cm} @{}}
    \toprule
    \textbf{Field} & \textbf{Required} & \textbf{Description} \\
    \midrule
    Name        & Yes & The public-facing title of the event. \\
    Description & No  & A summary of the event. \\
    Date \& Time & Yes & Date and start time in \texttt{yyyy-mm-dd hh:mm} format. \\
    Capacity    & No  & Maximum number of attendees. Leave blank for unlimited. \\
    Location    & Yes & The venue or address where the event is held. \\
    Price (USD) & No  & Ticket price. Leave empty or enter \texttt{0} for a
                        free event with no payment screen. \\
    Image URL   & No  & A direct URL to an externally hosted image used as
                        the event banner. \\
    Requires table signup & No & Enables table and seat selection for
                                  attendees. See Section~\ref{sec:table}. \\
    Requires bus signup   & No & Enables bus and seat selection for
                                  attendees. See Section~\ref{sec:bus}. \\
    \bottomrule
  \end{tabular}
\end{table}

\section{Table Signup Configuration}
\label{sec:table}

When \uielem{Requires table signup} is checked, you will be prompted to enter:
\begin{itemize}
  \item Number of tables at the venue.
  \item Seats per table.
\end{itemize}

Attendees will then choose their preferred table number and seat number
when registering. Selection is by number only --- there is no visual seating map.

\begin{notebox}
Table and seat assignments are visible in the Registered Attendees panel
on the event detail page.
\end{notebox}

\section{Bus Signup Configuration}
\label{sec:bus}

When \uielem{Requires bus signup} is checked, you will be prompted to enter:
\begin{itemize}
  \item Number of buses available.
  \item Seats per bus.
\end{itemize}

Attendees will choose their preferred bus number and seat number when registering.

\section{Viewing an Event}

Clicking an event from the \uielem{Events} list opens the event detail page.
What you see depends on whether your role matches the event's managing role:

\textbf{All staff} can see:
\begin{itemize}
  \item Event metadata: description, location, capacity, price.
  \item Seating configuration (tables and/or buses, if configured).
\end{itemize}

\textbf{Staff with the matching managing role} (and System Admins) additionally see:
\begin{itemize}
  \item \uielem{Edit event}, \uielem{Delete}, and
        \uielem{Download venue report} buttons.
  \item \uielem{Managing role} selector.
  \item \uielem{Registered Attendees} table with contact details, seat
        assignments, check-in status, and sign-up timestamps.
\end{itemize}

\section{Registered Attendees Table}

The Registered Attendees section shows a full table of everyone signed up,
with the following columns:

\begin{table}[H]
  \centering
  \caption{Registered Attendees Table Columns}
  \begin{tabular}{@{} p{3cm} p{10cm} @{}}
    \toprule
    \textbf{Column} & \textbf{Description} \\
    \midrule
    Attendee    & Full name and McMaster program. \\
    Contact     & Email address and phone number. \\
    Table Seat  & Assigned table and seat (if table signup enabled). \\
    Bus Seat    & Assigned bus and seat (if bus signup enabled). \\
    Check-in    & \texttt{Not checked in} or \texttt{Checked in}. \\
    Signed Up   & Date and time of registration. \\
    \bottomrule
  \end{tabular}
\end{table}

\section{Editing an Event}

\begin{enumerate}
  \item Open the event detail page.
  \item Click \uielem{Edit event}.
  \item Update the relevant fields.
  \item Click \uielem{Save} to apply changes.
\end{enumerate}

\section{Deleting an Event}

\begin{warnbox}
Deleting an event is \textbf{permanent and cannot be undone}. All associated
registrations and attendee data will also be removed.
\end{warnbox}

\begin{enumerate}
  \item Open the event detail page.
  \item Click \uielem{Delete}.
  \item Confirm the deletion in the dialog that appears.
\end{enumerate}


\chapter{Venue Reports}

\section{Downloading a Venue Report}

The venue report is a PDF summary of an event's details and registered
attendees, suitable for sharing with venue coordinators or transport providers.

\begin{enumerate}
  \item Navigate to \uielem{Events} and click the relevant event.
  \item Click \uielem{Download venue report} at the top of the page.
  \item The PDF will download automatically.
\end{enumerate}

\begin{notebox}
The \uielem{Download venue report} button is only visible to staff whose
role matches the event's managing role, and to System Admins.
\end{notebox}

\section{What the Venue Report Contains}

The report is a two-page PDF structured as a formal invoice:

\textbf{Page 1 --- Event Summary:}
\begin{itemize}
  \item Invoice number, issue date, and status.
  \item Event name, description, location, date/time, price per ticket,
        capacity, and prepared-by name.
  \item Seating configuration --- tables and seats per table.
  \item Transportation --- number of buses, seats per bus.
\end{itemize}

\textbf{Page 2 --- Registered Attendees:}
\begin{itemize}
  \item Total registered count.
  \item Table listing each attendee's name, program, email, table
        assignment, bus assignment, and sign-up date.
\end{itemize}

\begin{tipbox}
The venue report is ideal for sharing with venue coordinators, catering
teams, or transportation providers ahead of the event.
\end{tipbox}



\chapter{Notifications}


\section{Overview}

Staff can send push notifications to all registered attendees of an event
from the \uielem{Notifications} page. You can only send notifications for
events whose managing role matches one of your assigned roles. System Admins
can send notifications for any event.

\section{Sending a Notification}

\begin{enumerate}
  \item Navigate to \uielem{Notifications} in the sidebar.
  \item Use the \uielem{Event} dropdown to select the event. Only events
        you have permission to manage will appear.
  \item The form shows how many attendees will be notified.
  \item Optionally enter a \uielem{Message} for extra context.
  \item Click \uielem{Send Notification}.
\end{enumerate}

A confirmation banner will appear: \textit{``Notification sent to N
attendee(s).''} The button is then disabled for \textbf{60 minutes}.

\begin{warnbox}
There is a \textbf{60-minute cooldown} between notifications for the same
event. A warning banner will appear if a cooldown is active. Plan your
notifications carefully, especially close to event time.
\end{warnbox}

\begin{tipbox}
Only attendees \textbf{registered} for the selected event will receive
the notification.
\end{tipbox}



\chapter{Managing Users}


\section{Overview}

The \uielem{Users} page lists every registered MacSync user. Access level
determines what you can do on a user's profile:

\begin{table}[H]
  \centering
  \caption{Users Page --- Access by Privilege Level}
  \begin{tabular}{@{} p{3.5cm} p{9cm} @{}}
    \toprule
    \textbf{Privilege} & \textbf{What They Can See / Do} \\
    \midrule
    Staff (any role)  & View public profile: name, email, tickets purchased,
                        events attended, program, and ticket history. \\
    System Admin      & Everything above, plus Advanced Settings: assign/remove
                        roles, toggle System Admin, delete the user. \\
    \bottomrule
  \end{tabular}
\end{table}

\section{Viewing the Users List}

\begin{enumerate}
  \item Navigate to \uielem{Users} in the sidebar.
  \item The list shows all registered users with their name, email, and
        any special privilege badges (e.g.\ \texttt{SYSTEM\_ADMIN}).
  \item Use the \uielem{Search users by name or email} field to find a
        specific user.
\end{enumerate}

\section{Viewing a User Profile}

Click any user to open their profile, which displays:

\begin{itemize}
  \item Name, email, and active roles (shown top-right).
  \item Tickets purchased, events attended, and McMaster program.
  \item Full ticket history: event name, date, location, price, purchase
        date, seat assignment, and check-in status.
\end{itemize}

\section{Deleting a User (System Admin Only)}

\begin{warnbox}
Deleting a user is \textbf{permanent and cannot be undone}. All associated
data including tickets, registrations, and seat assignments will be removed.
\end{warnbox}

\begin{enumerate}
  \item Open the user's profile.
  \item Click \uielem{Show advanced settings}.
  \item Click \uielem{Delete user} and confirm.
\end{enumerate}



\chapter{Glossary}


\begin{description}[leftmargin=3.5cm, style=sameline]
  \item[Admin]         A built-in staff role. Can create and manage events
                       where their role matches the event's managing role.
  \item[Check-in]      The process of scanning a user's QR code at the
                       event entrance to mark them as attended.
  \item[Entry QR Code] A unique scannable code on a user's ticket used for
                       event check-in.
  \item[Event]         A scheduled activity created and managed within MacSync,
                       visible to all registered users on the mobile app.
  \item[Managing Role] The role assigned to an event that determines which
                       staff users can edit, delete, and manage it.
  \item[Member]        The default role for every registered MacSync user.
                       Can browse and register for events.
  \item[MES]           McMaster Engineering Society --- the organisation that
                       operates the MacSync platform.
  \item[RBAC]          Role-Based Access Control --- the permission system
                       governing what each user can do based on their roles.
  \item[System Admin]  A special privilege toggle granting unrestricted access
                       to all platform features, events, and settings.
  \item[Venue Report]  A two-page PDF summarising an event's details, seating,
                       transportation, and registered attendees.
\end{description}

\end{document}